% ===================== PARTE 11: VEREDICTO Y CIERRE (nueva) =====================
\parte{Veredicto y cierre}{Veredicto y cierre}
  {Las 15 afirmaciones del paper, una por una; los límites de esta réplica, sin rodeos; y lo que queda abierto.}{p11_portada}

\begin{frame}{El veredicto, afirmación por afirmación}
  \fuente{EXPLICACION.md, Parte 7; cifras de la tanda con Docker}
  \fontsize{6.4}{7.4}\selectfont\renewcommand{\arraystretch}{0.98}\setlength{\tabcolsep}{4pt}
  \begin{tabularx}{\linewidth}{@{}r p{0.33\linewidth} X l@{}}
    \toprule
    & \textbf{Lo que dice el paper} & \textbf{Resultado aquí} & \textbf{Veredicto}\\
    \midrule
    1 & Mucho más rápido que Libstdc++ & Gana con contención (0,06\por{} a 0,16\por); empata con $n = \pot{6}$ & \sirep\\
    2 & Mucho más rápido que Williams & Gana en los 6 escenarios (0,24\por{} a 0,73\por) & \sirep\\
    3 & Mucho más rápido que Folly & Gana 3, empata 2 y pierde 1; la brecha ya existe con 1 hilo & \norep\\
    4 & Con $u$ pequeño, supera a Anderson & +6\,\% con $n = 1$ y +65\,\% con $n = 10$ & \sirep\\
    5 & Con $n$ grande, supera a Anderson & Pierde hasta 18 hilos y empata desde 24 & \norep\\
    6 & Con mucha contención, Anderson gana & Anderson 1,87\por{} a 1,99\por & \sirep\\
    7 & \texttt{load}, siempre el de menor latencia & El más bajo en los tres casos; con contención, por poco & \sirep\\
    8 & \texttt{store}, segundo con $p = 1$ & Cuarto de cinco, por 5\,ns & \norep\\
    9 & \texttt{store}, el mejor con contención & Empate con Anderson & \enparte\\
    10 & \texttt{store \& reclaim} más lento con $p = 1$ & Los cinco & \sirep\\
    11 & Con contención, algunos más rápidos en \texttt{s\&r} & Sin Docker y 15 rondas, 1-4\,\%: ruido & \noconf\\
    12 & Williams, latencia casi siempre mayor & Sin contención sí; con contención, no siempre & \enparte\\
    13 & Anderson, poco sobrecoste; Autores, ninguno & Anderson hasta 469\,KB; los demás, bajo el límite & \sirep\\
    14 & La recarga reintenta como mucho $P$ veces & Máximo 8 con $P = 32$, en millones de recargas & \sirep\\
    15 & No hay \emph{use-after-free} & 0 errores con ASan y TSan, en lo probado & \sirep\\
    \bottomrule
  \end{tabularx}
  \nota{p11_veredicto}
\end{frame}

\begin{frame}{En una frase, y en un marcador}
  \fuente{EXPLICACION.md, Partes 0 y 7}
  \vspace{2mm}
  \begin{columns}[T]
    \column{0.24\textwidth}\cifra{\textcolor{verde}{9}}{se reproducen}
    \column{0.24\textwidth}\cifra{\textcolor{ambar}{2}}{en parte}
    \column{0.24\textwidth}\cifra{\textcolor{rojo}{3}}{no se reproducen}
    \column{0.24\textwidth}\cifra{\textcolor{gris}{1}}{no se confirma}
  \end{columns}
  \vspace{6mm}
  \caja{\textbf{La idea central del paper sale bien parada:} el algoritmo es correcto, cumple su cota, no acumula basura y lee rápido. Lo que no se sostiene aquí es la superioridad frente a Folly ni la ventaja con muchísimas pilas, y no se pudo reproducir la escala de 4 sockets.}
  \par\medskip
  \aviso{Sin Docker, con los mismos programas oficiales, la tabla sale igual salvo la fila 11, que deja de confirmarse con 15 rondas.}
  \nota{p11_frase}
\end{frame}

\begin{frame}{Limitaciones, dichas sin rodeos}
  \fuente{EXPLICACION.md, Parte 8}
  \begin{columns}[T]
    \column{0.32\textwidth}
    \sub{La máquina}
    \begin{itemize}\scriptsize
      \item No hay 4 sockets ni NUMA: se simulan; castigo de 1,35\por{} frente a unas 3\por.
      \item 24 núcleos, no 72, y sin SMT: la zona de más de 54 hilos queda fuera.
      \item Un solo módulo de RAM: limita $n = \pot{6}$.
      \item CPU híbrida (P y E) y frecuencia sin fijar.
    \end{itemize}
    \column{0.32\textwidth}
    \sub{El método}
    \begin{itemize}\scriptsize
      \item Recortes de 2\,s \por{} 5: desvío de hasta un 21\,\%, con sesgo.
      \item Valores del paper leídos de su PDF, no sus datos crudos.
      \item Dos medidas de Williams descartadas por errores nuestros.
      \item ASan y TSan cortos; Folly y Anderson sin TSan.
    \end{itemize}
    \column{0.32\textwidth}
    \sub{Las conclusiones}
    \begin{itemize}\scriptsize
      \item Folly, Williams entre clústeres y Autores con $n = \pot{6}$ y pocos hilos: \textbf{preguntas abiertas}.
      \item El techo de RAM se apoya en contadores, pero no está demostrado.
      \item Las causas no medidas se presentan como hipótesis.
    \end{itemize}
  \end{columns}
  \vfill
  \caja{Una prueba solo encuentra errores; no demuestra que no existan. La demostración de corrección es la del paper.}
  \nota{p11_limites}
\end{frame}

\begin{frame}{Lo que queda abierto, y para el equipo}
  \fuente{replica/completa/TRASPASO.md, «Lo que sigue abierto»; EXPLICACION.md, Parte 8}
  \begin{columns}[T]
    \column{0.5\textwidth}
    \sub{Preguntas abiertas}
    \begin{itemize}\footnotesize\setlength{\itemsep}{2pt}
      \item Por qué Folly rinde aquí mucho más que en el paper.
      \item Por qué Williams se hunde entre clústeres en el programa oficial, y más dentro de Docker ($-97$\,\%) que fuera ($-76$\,\%).
      \item Por qué Autores va por detrás con $n = \pot{6}$ y pocos hilos.
      \item Si la RAM es de verdad el techo: haría falta un segundo módulo.
      \item Por qué el recorte tiene sesgo.
    \end{itemize}
    \column{0.46\textwidth}
    \sub{Para el equipo}
    \begin{itemize}\footnotesize\setlength{\itemsep}{2pt}
      \item \texttt{bench.cpp}, línea 73: añadir \texttt{alignas(128)} a \texttt{Pila}.
      \item \texttt{replica/informe/02\_latencia.md}, línea 18: el paper pone 99,7\,\textmu s, no 86\,700\,ns.
      \item \texttt{ejecutar\_pruebas.sh}: a TSan le falta \texttt{setarch -R}.
      \item Quien active a Williams en latencia, que compile con \texttt{-mcx16}.
    \end{itemize}
  \end{columns}
  \nota{p11_abierto}
\end{frame}

{\setbeamercolor{background canvas}{bg=azul}\setbeamertemplate{footline}{}
\begin{frame}[plain]
  \vfill\centering
  {\color{white}\fontsize{28}{32}\selectfont\bfseries ¿Preguntas?\par}\bigskip
  {\color{azulclaro}\large David · Jheremy · Diego}\par\vfill
  {\color{white}\small\bfseries Código, datos y explicación completa}\par\medskip
  {\color{azulclaro}\small\repos}\par\medskip
  {\color{azulclaro}\scriptsize En la rama \texttt{replica-completa}: \texttt{replica/completa/EXPLICACION.md} y \texttt{replica/resultados/david-ultra9-285/}}
  \vspace*{6mm}
  \nota{p11_preguntas}
\end{frame}}
